\section{财产犯罪 - 盗窃罪}
\begin{enumerate}
  \item 和侵占罪一起考。
  \item \zd{他人占有（不要求他人所有）转为自己占有，对物平和}。区分：\zd{谁在占有}。
  \item 占有构成要件。和民法不一样。
        \begin{enumerate}
          \item 事实占有。实际控制范围内。用空间距离引诱你。\zd{第 1 步应先判断是不是主人在占有}。
                \begin{enumerate}
                  \item 控制范围：\zd{紧密 + 松弛占有}。
                  \item 占有的辅助者。\al{看仓库}。
                  \item 短暂遗忘。\al{去上卫生间，手机放在餐桌，隔桌拿走}。
                  \item 上下占有。\zd{授权大小，高度充分信赖}。
                        \begin{enumerate}
                          \item \al{舅舅让外甥帮卖货，外甥拿走了东西}。
                          \item \al{长途运输}。路线自己设计、没有押运员、自己一个人。
                        \end{enumerate}
                  \item 封缄物占有。\zd{外形物整体物占有；内容物拿走盗窃}。\al{保险柜、快递}。要有严格的防护保密措施。
                  \item 共同占有。\al{甲乙合资卖电脑，乙私自卖}。
                  \item 占有的转化。看\zd{管控力大小}（人员流动）。
                        \begin{enumerate}
                          \item \al{东西落在出租车，后面的乘客拿走}。转为出租车司机在占有。
                        \end{enumerate}
                \end{enumerate}
          \item 观念占有。\al{楼下停车，去北京出差}。
          \item 占有意思。
                \begin{enumerate}
                  \item \al{偷功德箱的钱}。寺庙在占有。
                  \item 死者占有。承认人死了，\zd{财物还有一定时间的延续性}。\al{甲杀仇人乙，发现竟然还有财物，拿走}。
                        \begin{enumerate}
                          \item 当场杀，杀人者拿。多数说，定盗窃。
                          \item 当场杀过后拿，一律定侵占。
                          \item 当场杀第三人拿 \yjt 侵占。
                          \item 家里杀家里拿 \yjt 盗窃。
                        \end{enumerate}
                \end{enumerate}
        \end{enumerate}
  \item 手段问题。\zd{对物平和}。\zd{不要求秘密性}。
        \begin{enumerate}
          \item \al{甲在公交车上扒窃乙的钱包，乙浑然不知，但周围乘客都看到了不吭声}。
          \item \al{小偷入户盗窃，老人被惊醒不敢吭声}。
          \item \al{甲以为李是住宅主人，打晕然后进去拿走东西，但主人其实是王某}。
        \end{enumerate}
  \item 特殊盗窃类型。\zd{数额标准降低}，不要求较大（普通盗窃 2000 左右），只要是刑法值得保护的财物。
        \begin{enumerate}
          \item 多次；携带凶器；扒窃。
          \item 入\zd{户}盗窃（不是处罚情节）。\al{煎饼摊盗窃}。有人居住的生活场所。\zd{不要和入户抢劫混了}。
        \end{enumerate}
\end{enumerate}
